The phase convention is as follows, assuming a positive rf voltage:
\verb|PHASE=90| is the crest for acceleration.  \verb|PHASE=180| is the stable
phase for a storage ring above transition without energy losses.

{\bf Body focus models}: By default, no body focusing (i.e., focusing along the
cavity) is included. Two options are included:
\begin{itemize}
\item[SRS] --- Simplified Rosenzweig-Serafini, based on Rosenzweig and Serafini, Phys. Rev. E 49 (2),
1599.  As suggested by N. Towne (NSLS), I simplified this to assume a pure pi-mode
standing wave. To use this, one must also set \verb|STANDING_WAVE=1|.
\item[TW1] --- Traveling Wave 1, which includes the radial electric and azimuthal
magnetic field components for a pure traveling wave given by 
\begin{equation}
        E_z = E_0 \sin \left(\omega t - k z + \phi\right)
\end{equation}
To use this, one must set \verb|STANDING_WAVE=0| (the default). One must also have \verb|N_KICKS|$\geq 10$.
\end{itemize}

The \verb|CHANGE_T| parameter may be needed for reasons that stem from
{\tt elegant}'s internal use of the total time-of-flight as the
longitudinal coordinate.  If the accelerator is very long or a large
number of turns are being tracked, rounding error may affect the
simulation, introducing spurious phase jumps.  By setting
\verb|CHANGE_T=1|, you can force {\tt elegant} to modify the time
coordinates of the particles to subtract off $N T_{rf}$, where
$T_{tf}$ is the rf period and $N = \lfloor t/T_{tf}+0.5\rfloor$.  If
you are tracking a ring with rf at some harmonic $h$ of the revolution
frequency, this will result in the time coordinates being relative to
the ideal revolution period, $T_{rf}*h$.  If you have multiple rf
cavities in a ring, you need only use this feature on one of them.
Also, you can use \verb|CHANGE_T=1| if you simply prefer to have the
offset time coordinates in output files and analysis. 

As of version 2026.4, if there are rf cavities that are not harmonically-related
to the cavity with \verb|CHANGE_T=1|, there will be phase adjustments made to
correct for the time offsets. The time offsets are accumulated in extended-precision
registers, then used to compute the phase adjustments modulo $2\pi$, which reduces
phase jumps due to insufficient precision.
In addition, elements with longer-term, non-resonant time dependence (e.g., \verb|BUMPER|
or a damped \verb|RFMODE|) will behave properly as they'll be provided with the accumulated
time offset.
